\documentclass[12pt]{article}
\usepackage[utf8]{inputenc}
\usepackage[T1]{fontenc}
\usepackage{hyperref}
\title{Totomachapan Zapotec Language Resources for K-12 Curriculum Development}
\author{Compiled by Claude Sonnet 5.5 (claude-sonnet-5-5) for Andr\'as Kornai}
\date{}
\begin{document}
\maketitle

\section{Introduction}
Totomachapan Zapotec (ISO 639-3 zph; Glottocode toto1308), also called Western Zimatl\'an Zapotec, is an endangered Otomanguean language of Oaxaca, Mexico. Sources describe it as spoken mainly by older adults and as unwritten. Vitality ratings differ between sources. Very little material exists. The search found no primer, textbook, dictionary, grammar or corpus.

\section{Primary Educational Resources}
No primers, textbooks or school programmes were found in the search results. The language is not known to be taught in schools.

\section{Linguistic and Reference Materials}
\begin{itemize}
\item Wikipedia stub article on the language: \url{https://en.wikipedia.org/wiki/Totomachapan_Zapotec}
\item Wikitongues language page: \url{https://wikitongues.org/languages/zph/}
\item Verbix Endangered Languages entry: \url{https://docs.verbix.com/EndangeredLanguages/Totomachapam}
\end{itemize}

\section{Religious and Community Resources}
No community organisation or scripture translation was found. Global Recordings Network (below) offers Bible-teaching audio.

\section{Digital and Audio Resources}
\begin{itemize}
\item Global Recordings Network, Zapotec, Totomachapan (evangelistic and Bible-teaching audio, slideshows): \url{https://www.globalrecordings.net/en/language/zph}
\end{itemize}

\section{Curriculum Development Context}
The language is unwritten and has few elderly speakers, and no teaching materials were found. A curriculum would first need new documentation work (recordings, a practical orthography, a lexicon) with the community. This report does not claim to cover all existing materials.
\end{document}